一个质量为 \(m>0\)、电荷为 \(e\ne0\) 的非相对论离子在真空平面四极电势
\[
\Phi(x,y,t)=\frac{U+V\cos\Omega t}{2r_0^2}(x^2-y^2),\quad \Omega,r_0>0
\]
中运动；理想模型忽略磁力、空间电荷与电极端部场。
(1) 验证真空Laplace方程；将粒子的 \(x,y\) 运动写成 \(z''+(a_z-2q_z\cos2\tau)z=0\)、\(\tau=\Omega t/2\)，解释为何只有直流鞍点时不能双轴束缚。
(2) 令 \(U=0,\ \epsilon=eV/(mr_0^2)\)，且 \(|\epsilon|\ll\Omega^2\)。用快慢尺度求微运动、世俗振动频率及两维赝势，要求给出适用的小参数与量纲检验。
(3) 对近乎静止的慢位置 \(X,Y\)，分别求微运动一个射频周期的平均动能与射频电势能；由平均Lagrangian严格辨认有效赝势，给出原系统显含时间的能量变化率。
(4) 加入弱直流 \(U\) 后求双轴世俗频率平方、近似稳定域及其临界退化，讨论 \(V\to0\) 与固定 \(U,V\) 的 \(\Omega\to\infty\) 极限，说明不能以近似平均理论取代精确Mathieu稳定图。